\documentclass[11pt]{article}
\usepackage[margin=1in]{geometry}
\usepackage{amsmath,amssymb,amsthm,mathtools}
\usepackage[hidelinks]{hyperref}
\usepackage{enumitem}
\usepackage{microtype}

\newtheorem{theorem}{Theorem}
\newtheorem{lemma}[theorem]{Lemma}
\newtheorem{proposition}[theorem]{Proposition}
\newtheorem{remark}[theorem]{Remark}

\DeclareMathOperator{\diam}{diam}
\newcommand{\N}{\mathbb N}
\newcommand{\eps}{\varepsilon}
\newcommand{\ph}{\varphi}
\DeclareMathOperator{\supp}{supp}

\title{A Converse Estimate for the Schlumprecht-Space Szlenk Radius}
\author{Draft proof note}
\date{June 2026}

\begin{document}
\maketitle

\begin{abstract}
Let $S$ be Schlumprecht's space, with
\[
 \|x\|_S=\max\left\{\|x\|_\infty,
 \sup_{E_1<\cdots<E_n}\frac1{\ph(n)}\sum_{i=1}^n\|E_i x\|_S\right\},
 \qquad \ph(n)=\log_2(n+1).
\]
Kochanek and Miarka proved [KM]
\[
 r_{S^*}(\eps)\le \inf_{n\in\N}
 \frac{\log_2(n+2)-\eps/2}{\log_2(n+1)},\qquad 0<\eps<2,
\]
and asked whether the converse inequality holds.  This note gives a proof of the converse.
The central point is a strengthened scalar majorant estimate.  It replaces a weaker estimate
that is not sufficient for the local recursive majorant argument.
\end{abstract}

\section{The logarithmic product estimate}

Throughout, put
\[
 \ph(n)=\log_2(n+1),\qquad
 B_i=\ph(i+1),\qquad
 \Delta_i=\ph(i+1)-\ph(i) \quad (i\ge 1).
\]

\begin{lemma}[monotonicity]
The function
\[
 h(u)=\frac{u\log u}{u-1}\qquad (u>1)
\]
is increasing.
\end{lemma}

\begin{proof}
A direct differentiation gives
\[
 h'(u)=\frac{u-1-\log u}{(u-1)^2}\ge 0,
\]
because $\log u\le u-1$ for $u>0$.
\end{proof}

\begin{lemma}[product lemma]
For all $i,j\ge 1$, there exists $\ell\in\N$, namely
\[
 \ell=(i+1)(j+1)-1,
\]
such that
\[
 B_\ell\le B_iB_j,
 \qquad
 \Delta_\ell\le \Delta_i\Delta_j.
\]
Equivalently, for all $p,q\ge 1$,
\begin{align}
 \log_2(pq+1)&\le \log_2(p+1)\log_2(q+1), \label{eq:prod-B}\\
 \log_2\left(1+\frac1{pq}\right)&\le
 \log_2\left(1+\frac1p\right)\log_2\left(1+\frac1q\right). \label{eq:prod-Delta}
\end{align}
\end{lemma}

\begin{proof}
We prove the two logarithmic inequalities.  For fixed $p\ge 1$, define
\[
 F_p(q)=\frac{\log(pq+1)}{\log(q+1)}\qquad(q>1),
\]
where $\log$ is the natural logarithm.  Its derivative has the sign of
\[
 p(q+1)\log(q+1)-(pq+1)\log(pq+1).
\]
Writing $u=q+1$ and $v=pq+1=1+p(u-1)$, this sign is non-positive because
\[
 \frac{u\log u}{u-1}\le \frac{v\log v}{v-1}
\]
by the preceding monotonicity lemma.  Hence $F_p$ is decreasing, and therefore
\[
 \frac{\log(pq+1)}{\log(q+1)}\le
 \frac{\log(p+1)}{\log 2}.
\]
This is exactly \eqref{eq:prod-B} after division by $\log 2$.

For \eqref{eq:prod-Delta}, put $a=1/p$ and $b=1/q$, so $0<a,b\le 1$.  For fixed $a$ define
\[
 G_a(b)=\frac{\log(1+ab)}{\log(1+b)}\qquad(0<b\le 1).
\]
The derivative has the sign of
\[
 a(1+b)\log(1+b)-(1+ab)\log(1+ab).
\]
Writing $u=1+b$ and $v=1+ab=1+a(u-1)$, this sign is non-negative, again by the monotonicity of
$u\log u/(u-1)$.  Thus $G_a$ is increasing on $(0,1]$, and
\[
 \frac{\log(1+ab)}{\log(1+b)}\le
 \frac{\log(1+a)}{\log 2}.
\]
This is \eqref{eq:prod-Delta}.

Now take $p=i+1$ and $q=j+1$.  If $\ell=pq-1$, then
\[
 B_\ell=\ph(\ell+1)=\log_2(pq+1),
 \qquad
 B_i=\log_2(p+1),\quad B_j=\log_2(q+1),
\]
and
\[
 \Delta_\ell=\log_2\left(1+\frac1{pq}\right),
 \quad
 \Delta_i=\log_2\left(1+\frac1p\right),
 \quad
 \Delta_j=\log_2\left(1+\frac1q\right).
\]
The two desired inequalities follow.
\end{proof}

\section{The scalar majorant}

For $0<\alpha<1$ set
\[
 a_\alpha=\inf_{n\in\N}\frac{\ph(n+1)-\alpha}{\ph(n)}.
\]
Equivalently,
\begin{equation}
 \ph(n)a_\alpha+\alpha\le \ph(n+1)\qquad(n\in\N).
 \label{eq:a-alpha-def}
\end{equation}
For $u\ge 0$ define
\[
 M_\alpha(u)=\max\left\{u,\alpha,
 \sup_{m\in\N}\frac{\ph(m)u+\alpha}{\ph(m+1)}\right\}.
\]
In particular, if $0\le u\le a_\alpha$, then $M_\alpha(u)\le 1$, by
\eqref{eq:a-alpha-def} and $a_\alpha\le 1$.

\begin{lemma}[strengthened scalar estimate]
Let $0<\alpha<1$, $0\le s\le t$, and $k\ge 2$.  Then
\begin{equation}
 M_\alpha(s)+\min\{\ph(k-1)t,\ph(k)t-s\}
 \le \ph(k)M_\alpha(t).
 \label{eq:strong-scalar}
\end{equation}
Consequently, the same inequality holds for $0\le s\le t\le a_\alpha$.
\end{lemma}

\begin{proof}
Let
\[
 a=\ph(k-1),\qquad b=\ph(k),\qquad \delta=b-a.
\]
It suffices to verify \eqref{eq:strong-scalar} for each affine term in the maximum defining
$M_\alpha(s)$.

For the term $s$, use the second member of the minimum:
\[
 s+\min\{at,bt-s\}\le s+(bt-s)=bt\le bM_\alpha(t).
\]
For the term $\alpha$, use the first member of the minimum and the $(k-1)$-st affine term in
$M_\alpha(t)$:
\[
 \alpha+\min\{at,bt-s\}\le \alpha+at
 \le bM_\alpha(t).
\]

It remains to consider the $m$-th affine term.  Put
\[
 c=\ph(m),\qquad d=\ph(m+1),\qquad \eta=d-c.
\]
Since $0<\delta<1$, for fixed $t$, the function
\[
 s\mapsto \frac{cs+\alpha}{d}+\min\{at,bt-s\}
\]
is increasing on $[0,\delta t]$ and decreasing on $[\delta t,t]$.  Hence its maximum over
$0\le s\le t$ occurs at $s=\delta t$, where it equals
\[
 \frac{\alpha+t(c\delta+ad)}{d}
 =\frac{\alpha+t(bd-\delta\eta)}{d}.
\]
Thus it remains to prove
\begin{equation}
 M_\alpha(t)\ge
 \frac{\alpha+t(bd-\delta\eta)}{bd}.
 \label{eq:scalar-reduced}
\end{equation}
Apply the product lemma with $i=k-1$ and $j=m$.  There is $\ell\in\N$ such that, writing
$D=B_\ell=\ph(\ell+1)$ and $\lambda=\Delta_\ell=\ph(\ell+1)-\ph(\ell)$,
\[
 D\le bd,
 \qquad
 \lambda\le \delta\eta.
\]
The $\ell$-th affine term in $M_\alpha(t)$ gives
\[
 M_\alpha(t)
 \ge \frac{\ph(\ell)t+\alpha}{\ph(\ell+1)}
 =t+\frac{\alpha-\lambda t}{D}.
\]
If $\alpha-\delta\eta t<0$, then the right-hand side of \eqref{eq:scalar-reduced} is below
$t$, while $M_\alpha(t)\ge t$.  If $\alpha-\delta\eta t\ge 0$, then $\alpha-\lambda t\ge
\alpha-\delta\eta t\ge 0$ and $D\le bd$, so
\[
 M_\alpha(t)
 \ge t+\frac{\alpha-\delta\eta t}{bd}
 =\frac{\alpha+t(bd-\delta\eta)}{bd}.
\]
This proves \eqref{eq:scalar-reduced}, and hence \eqref{eq:strong-scalar}.
\end{proof}

\begin{remark}
The scalar estimate needed in the recursive norm argument is \eqref{eq:strong-scalar}, not merely
\[
 M_\alpha(s)+\min\{\ph(k-1)t,\ph(k)t-s\}\le \ph(k).
\]
The extra factor $M_\alpha(t)$ on the right is the local normalization required when one estimates
the norm inside a restriction $F$.
\end{remark}

\section{The far-coordinate perturbation lemma}

We shall use the standard iterative construction of the Schlumprecht norm.  Set
$\|x\|_0=\|x\|_\infty$ and
\[
 \|x\|_{r+1}=\max\left\{\|x\|_\infty,
 \sup_{k\ge 2}\sup_{E_1<\cdots<E_k}
 \frac1{\ph(k)}\sum_{i=1}^k\|E_i x\|_r\right\}.
\]
Then $\|x\|_S=\sup_r\|x\|_r$.  The case $k=1$ is redundant, since $\ph(1)=1$.

\begin{lemma}[comparison by local majorants]
Let $y\in c_{00}$.  Suppose $\rho(F)\ge 0$ is assigned to every finite $F\subset\N$ and satisfies
\begin{align*}
 \|Fy\|_\infty&\le \rho(F),\\
 G\subset F&\Longrightarrow \rho(G)\le \rho(F),\\
 \sum_{i=1}^k\rho(E_i)&\le \ph(k)\rho(F)
\end{align*}
whenever $k\ge 2$ and $E_1<\cdots<E_k$ are finite subsets of $F$.  Then
\[
 \|Fy\|_S\le \rho(F)
 \qquad(F\subset\N\text{ finite}).
\]
\end{lemma}

\begin{proof}
Induct on the iterative norms.  The estimate is true for $\|\cdot\|_0$ by the first assumption.
Assume it is true for $\|\cdot\|_r$, and fix a finite set $F$ and a collection
$E_1<\cdots<E_k$, $k\ge 2$.  Put $G_i=E_i\cap F$ and discard the empty intersections.
If no $G_i$ remains, the corresponding sum is zero.  If exactly one set $G$ remains, then
\[
 \|Gy\|_r\le \rho(G)\le \rho(F)\le \ph(k)\rho(F).
\]
If $q\ge 2$ sets $G_1<\cdots<G_q$ remain, then, since $\ph(q)\le \ph(k)$,
\[
 \sum_{j=1}^q\|G_jy\|_r
 \le \sum_{j=1}^q\rho(G_j)
 \le \ph(q)\rho(F)
 \le \ph(k)\rho(F).
\]
Thus every admissible sum occurring in the definition of $\|Fy\|_{r+1}$ is bounded by
$\ph(k)\rho(F)$, and so $\|Fy\|_{r+1}\le \rho(F)$.  Taking the supremum over $r$ proves the
claim.
\end{proof}

\begin{lemma}[far-coordinate unit-ball lemma]
Let $0<\alpha<1$.  If $x\in c_{00}$ and $\|x\|_S<a_\alpha$, then, for every
$N>\max\supp x$,
\[
 \|x+\alpha e_N\|_S\le 1,
 \qquad
 \|x-\alpha e_N\|_S\le 1.
\]
\end{lemma}

\begin{proof}
By $1$-unconditionality, it is enough to prove the assertion for $x\ge 0$ and
$y=x+\alpha e_N$.  For each finite $F\subset\N$ define
\[
 \rho(F)=
 \begin{cases}
  \|Fx\|_S, & N\notin F,\\
  M_\alpha(\|Fx\|_S), & N\in F.
 \end{cases}
\]
Since $\|Fx\|_S\le \|x\|_S<a_\alpha$, we have $\rho(F)\le 1$ whenever $N\in F$.
Also $\|Fy\|_\infty\le \rho(F)$: this is immediate if $N\notin F$, and if $N\in F$ it follows
from $M_\alpha(u)\ge u$ and $M_\alpha(u)\ge \alpha$.  The map $M_\alpha$ is increasing, being
the maximum of increasing affine functions; therefore $G\subset F$ implies $\rho(G)\le\rho(F)$.

We verify the local majorant inequality.  Fix $F$ and $E_1<\cdots<E_k\subset F$, $k\ge 2$.
If $N\notin F$, then
\[
 \sum_i\rho(E_i)=\sum_i\|E_i x\|_S\le \ph(k)\|Fx\|_S=\ph(k)\rho(F).
\]
Assume now that $N\in F$, and put $t=\|Fx\|_S$.

If none of the $E_i$ contains $N$, then
\[
 \sum_i\rho(E_i)=\sum_i\|E_i x\|_S\le \ph(k)t\le \ph(k)M_\alpha(t)=\ph(k)\rho(F).
\]
If exactly one of them, say $E_j$, contains $N$, put $s=\|E_jx\|_S$.  The pieces not meeting
$N$ satisfy the two bounds
\[
 \sum_{i\ne j}\|E_i x\|_S\le \ph(k-1)t,
 \qquad
 \sum_{i\ne j}\|E_i x\|_S\le \ph(k)t-s.
\]
The first uses the $k-1$ pieces not meeting $N$; the second uses the full $k$-tuple, including
$E_jx$.  Therefore the strengthened scalar estimate gives
\[
 \sum_i\rho(E_i)
 \le M_\alpha(s)+\min\{\ph(k-1)t,\ph(k)t-s\}
 \le \ph(k)M_\alpha(t)
 =\ph(k)\rho(F).
\]
The comparison lemma now yields $\|Fy\|_S\le\rho(F)$ for all finite $F$.  Taking
$F=\supp x\cup\{N\}$ gives
\[
 \|x+\alpha e_N\|_S\le M_\alpha(\|x\|_S)\le 1.
\]
The estimate for $x-\alpha e_N$ follows from unconditionality.
\end{proof}

\section{The Szlenk-radius formula}

Recall that for a Banach space $X$,
\[
 r_X(\eps)=\sup\{r>0:rB_{X^*}\subset s_\eps B_{X^*}\}.
\]
For $X=S^*$ we identify $X^*=S^{**}$ with $S$.

\begin{theorem}
For every $0<\eps<2$,
\[
 r_{S^*}(\eps)=
 \inf_{n\in\N}\frac{\log_2(n+2)-\eps/2}{\log_2(n+1)}.
\]
\end{theorem}

\begin{proof}
Let $\beta=\eps/2$ and
\[
 a_\beta=\inf_{n\in\N}\frac{\ph(n+1)-\beta}{\ph(n)}.
\]
Kochanek and Miarka proved [KM] $r_{S^*}(\eps)\le a_\beta$.  We prove the reverse inequality.

First observe that $\alpha\mapsto a_\alpha$ is continuous on $(0,1)$; indeed, each function
\[
 \alpha\mapsto \frac{\ph(n+1)-\alpha}{\ph(n)}
\]
is $1$-Lipschitz because $\ph(n)\ge 1$, and the infimum of a family of $1$-Lipschitz functions
is $1$-Lipschitz.  Hence, if $x\in c_{00}$ and $\|x\|_S<a_\beta$, we can choose
$\alpha\in(\beta,1)$ such that $\|x\|_S<a_\alpha$.

For every $N>\max\supp x$, the far-coordinate lemma gives
\[
 x+\alpha e_N\in B_S,
 \qquad
 x-\alpha e_N\in B_S.
\]
Since $S$ is reflexive, the basis vector sequence $(e_N)$ is weakly null; for $X=S^*$ this is also
weak-star null in $X^*=S$.  Thus $x\pm\alpha e_N$ converge weak-star to $x$.  Moreover,
\[
 \|(x+\alpha e_N)-(x-\alpha e_N)\|_S=2\alpha>\eps.
\]
Therefore every weak-star neighborhood of $x$ intersects $B_S$ in a set of diameter greater than
$\eps$, and so $x\in s_\eps B_S$.

Now fix $0<r<a_\beta$.  If $z\in rB_S$, choose $z_m\in c_{00}$ with $z_m\to z$ in norm.
For all large $m$, $\|z_m\|_S<a_\beta$, hence $z_m\in s_\eps B_S$.  The set $s_\eps B_S$ is
weak-star closed, hence norm closed: its complement in $B_S$ is relatively weak-star open by the
definition of the Szlenk derivation.  Consequently $z\in s_\eps B_S$.  Thus
$rB_S\subset s_\eps B_S$ for every $r<a_\beta$, and so $r_{S^*}(\eps)\ge a_\beta$.
Combining this with the known upper estimate gives the formula.
\end{proof}

\section{Check against the conditional packet}

The conditional packet identifies the correct perturbative mechanism, but the majorant step needs
the local normalized estimate
\[
 \sum_i\rho(E_i)\le \ph(k)\rho(F),
\]
not merely a global estimate by $\ph(k)$.  The strengthened scalar estimate above is exactly the
missing normalized statement.  After this replacement, the far-coordinate lemma and the Szlenk
radius converse go through as written above.

\begin{thebibliography}{9}
\bibitem{KM} T. Kochanek and M. Miarka, \emph{When is the Szlenk derivation of a dual unit ball another ball?}, arXiv:2409.05516 (2024).
\bibitem{S} T. Schlumprecht, \emph{An arbitrarily distortable Banach space}, Israel J. Math. 76 (1991), 81--95.
\end{thebibliography}

\end{document}
